\chapter{The Bershadsky--Polyakov algebra}\label{chap:bershadsky-polyakov}\label{ch:bershadsky-polyakov}

The Bershadsky--Polyakov algebra is the first non-principal
Drinfeld--Sokolov reduction in which the bar--cobar complex, the
Verdier companion, and the shadow tower separate visibly.  The
self-transpose partition $(2,1)$ produces a same-family companion at
level $-k-6$ under the finite-type-or-completed Verdier package and the
subregular DS/bar transport theorem.  On the regular shadow locus of
Convention~\ref{conv:bp-generic-loci}, the Virasoro $T$-line has
class~$\mathbf{M}$ and the residual-current $J$-line has
class~$\mathbf{G}$; the global class is their channel supremum.

The central-charge identity $K^c_{\mathcal{B}} \equiv 50$
(Theorem~\ref{thm:bp-koszul-conductor-polynomial}) holds as a
polynomial identity in $\mathbb{Q}(k)$.  Fehily--Kawasetsu--Ridout
\cite[Definition~2.1 and Eq.~\textup{(}2.2\textup{)}]{FKR20} define an
ordinary bosonic vertex algebra with even strong generators
$J,G^+,G^-,T$.  Consequently the modular characteristic requires its
own genus-$1$ curvature computation.  The current scalar ledger is
\[
K^c(\mathrm{BP}_k,\mathrm{BP}_{-k-6})=50,\qquad
\Delta^{\mathrm{DS}}_{\mathrm{BP}}(k)
:=c_{\mathrm{BP}}(k)-c_{\mathrm{aff}}(k)=-(6k+1).
\]
The mixed-class profile
$(T\colon\mathbf{M},\;J\colon\mathbf{G},\;G^{\pm}\colon\text{even charged})$
is a non-principal-DS instance of the channel-maximum rule of
Chapter~\ref{chap:shadow-quadrichotomy-platonic}: the global
shadow class is $\mathbf{M}$.

\begin{table}[ht]
\centering
\small
\caption{Theorems A, B, C, D, H for the Bershadsky--Polyakov
algebra $\mathcal{B}^k = \mathcal{W}^k(\mathfrak{sl}_3, f_{(2,1)})$
on the BP working loci of Convention~\ref{conv:bp-generic-loci}.}
\label{tab:bp-five-theorems}
\begin{tabular}{lll}
\toprule
\textbf{Theorem} & \textbf{Statement for $\mathcal{B}^k$}
 & \textbf{Hypothesis / reference} \\
\midrule
A (reconstruction) &
 $\epsilon_{\mathcal B^k}\colon
 \Omega_X^{\mathrm{cont}}\bar B_X(\mathcal B^k)\to\mathcal B^k$
 in the raw finite-window or completed pro-conilpotent ambient
 & Thm~\ref{thm:bar-cobar-inversion-qi}; package $H_1$ \\
B (quadratic recognition) &
 $q_{\mathcal B^k}\colon(\mathcal B^k)^{\mathrm i}
 \to\bar B_X(\mathcal B^k)$ on the BP quadratic Koszul locus
 & comparison package $H_{\mathrm{CL}}$ \\
C (complementarity) &
 $K^c_{\mathcal{B}}=50$ exactly; the modular-characteristic
 complementarity is open
 & Prop~\ref{prop:bp-complementarity} \\
D (modular char.) &
 $\kappa(\mathcal{B}^k)$ is the genus-$1$ curvature coefficient;
 its closed formula is open
 & Prop~\ref{prop:bp-kappa} \\
H (chiral Hochschild) &
$H_H(\mathcal B^k;S_{\mathrm{BP}})$ gives
$\operatorname{Supp}\ChirHoch_X^\bullet(\mathcal B^k)
\subseteq S_{\mathrm{BP}}$; the polynomial $1+t^2$ follows after
$S_{\mathrm{BP}}\subseteq\{0,1,2\}$, the degree-two pairing, and the
fixed-product BP derivation calculation
 & Thm~\ref{thm:hochschild-polynomial-growth};
 Prop~\ref{prop:bp-curve-level-chirhoch} \\
\bottomrule
\end{tabular}
\end{table}

\begin{table}[ht]
\centering
\small
\caption{Shadow archetype data for the Bershadsky--Polyakov
algebra.}\label{tab:bp-shadow-archetype}
\begin{tabular}{ll}
\toprule
\textbf{Invariant} & \textbf{Value} \\
\midrule
Archetype tag & $\mathbf{M}$ (mixed-channel: $T$-channel $\mathbf{M}$,
 $J$-channel $\mathbf{G}$, $G^\pm$ even charged fields);
 global class is the channel supremum \\
Shadow depth $r_{\max}$ & $\infty$ (on $T$-line);
 $r_{\max}^J = 2$ (on $J$-line) \\
$\kappa(\mathcal{B}^k)$ &
 open genus-$1$ curvature computation \\
Central charge & $c = -(2k{+}3)(3k{+}1)/(k{+}3)$
 (Kac--Roan--Wakimoto minimal reduction) \\
Central-charge conductor & $K^c_{\mathcal{B}} = c(k) + c(-k{-}6) = 50$ \\
Modular complementarity & open; central-charge complementarity equals $50$ \\
Central midpoint & $c = 25$ on the complex fibre $k = -3 \pm 2i$ \\
Critical boundary & $k = -3$, the common pole \\
$T$-line poles & $c=0$ or $c=-22/5$ \\
$J/G$ degeneracies & $k=-3/2$ for $J_{(1)}J$;
 $k=-1,-3/2$ for $G^+_{(2)}G^-$ \\
$J$-line level & $(2k+3)/3$ in Feigin--Semikhatov normal form \\
$r$-matrix $r(z)$ & Multi-channel:
 $T$-channel has poles at $z^{-3}, z^{-1}$;
 $J$-channel has pole at $z^{-1}$ \\
Partition & $(2,1)$ (self-transpose, hook-type) \\
DS parent & $\widehat{\mathfrak{sl}}_3$ at the minimal nilpotent \\
\bottomrule
\end{tabular}
\end{table}

\begin{remark}[$5\times5$ $\kappa$-stratification row for $\mathcal{B}^k$]
\label{rem:bp-5x5-kappa-row}
\index{$\kappa$-stratification matrix!Bershadsky--Polyakov row}
We work in the Open quadrant at Beilinson level~$1$, with
$A_b=\mathcal{B}^k$ in the Feigin--Semikhatov chart generated by
$(T,J,G^+,G^-)$, and with
$k\in U_{\mathrm{OPE}}\cap U_{\mathrm{nd}}\cap U_T$.

On the standard Open-quadrant Beilinson tower, the five
$\kappa$-measurements
$(\kappa_{\mathrm{cat}},\kappa^{\mathrm{Hodge}}_{\mathrm{ch}},
\kappa^{\mathrm{Heis}}_{\mathrm{ch}},\kappa_{\mathrm{BKM}},
\kappa_{\mathrm{fiber}})$ specialise on $\mathcal{B}^k$ to
\[
\bigl(\kappa_{\mathrm{cat}}(\mathcal{B}^k),\;
\kappa^{\mathrm{Hodge}}_{\mathrm{ch}}(\mathcal{B}^k),\;
\kappa^{\mathrm{Heis}}_{\mathrm{ch}}(\mathcal{B}^k),\;
\kappa_{\mathrm{BKM}}(\mathcal{B}^k),\;
\kappa_{\mathrm{fiber}}(\mathcal{B}^k)\bigr)
=\bigl(0,\;\kappa_{\mathrm{BP}}^{g=1}(k),\;
(2k{+}3)/3,\;\text{unassigned},\;\text{unassigned}\bigr),
\]
where $\kappa_{\mathrm{BP}}^{g=1}(k)$ is
\ClaimStatusOpen{} pending the full genus-$1$ curvature computation.
The Heisenberg slot $(2k+3)/3$ is the residual-current level
$k_J=J_{(1)}J$ of the Feigin--Semikhatov normal form,
i.e.\ the $J$-line scalar $S_2^J=\kappa_J$ of
Proposition~\ref{prop:bp-jline-depth}; the Hodge slot is the
genus-$1$ scalar defined in Proposition~\ref{prop:bp-kappa}; the categorical slot is set to
$0$ by convention in the absence of a CY host
(Table~\ref{tab:five-by-five-kappa-stratification},
Footnote *).
The two right-hand entries await a Borcherds denominator and a compact
CY fibre attached to $\mathcal{B}^k$ as a chart algebra. The five
entries are five distinct invariants.  Their comparison proceeds through the
averaging morphism
$\Av:\fg^{E_1}\to\fg^{\mathrm{mod}}$. An injection of
$\mathcal{B}^k$ into a row of the
$5\times5$ $\kappa$-stratification matrix endowed with all five
columns would require a companion CY${}_d$ datum supplied by
Vol~III through the two-stage factorisation
$\Phi_d^{(\Sigma_{d-1},C)}=\mathrm{Sp}^{\mathrm{ch}}_{\Sigma_{d-1},C}
\circ\Phi_d^{\mathrm{FA}}$ at a named pair $(\Sigma_{d-1},C)$
and receives its target from a named pair $(\Sigma_{d-1},C)$.  The
construction of such a pair for $\mathcal{B}^k$ is an open Vol.~III
problem.
\end{remark}

\begin{remark}[Archetype tag: BP is mixed-channel class~$\mathbf{M}$]
\label{rem:bp-archetype-supremum}
\index{Bershadsky--Polyakov algebra!archetype tag}
\index{five-archetype dichotomy!channel-supremum rule}
We work on the Open-quadrant chart $A_b=\mathcal{B}^k$ at
Beilinson level~$1$, with
$k\in U_{\mathrm{nd}}\cap U_T$, and use the channel-supremum rule of
Chapter~\ref{chap:shadow-quadrichotomy-platonic}.

The global archetype tag of $\mathcal{B}^k$ is class~$\mathbf{M}$
(``mixed channel'') in the G/L/C/M/B taxonomy.
The channel-supremum rule
(Chapter~\ref{chap:shadow-quadrichotomy-platonic}) reads the global
$r_{\max}$ as the maximum of per-line $r_{\max}$, here
$\max(\infty,2)=\infty$, hence class $\mathbf{M}$. The
class-$\mathbf{B}$ row is reserved for Mukai-enhanced K3-type rows
with a Borcherds denominator (Vol~III's
$\mathcal H_{\mathrm{Muk}}(K3)$ row at
$(\kappa_{\mathrm{cat}},\kappa^{\mathrm{Hodge}}_{\mathrm{ch}},
\kappa^{\mathrm{Heis}}_{\mathrm{ch}},\kappa_{\mathrm{BKM}},
\kappa_{\mathrm{fiber}})=(0,0,3,5,24)$).  BP carries the Open-quadrant
chart data described above; the $\Delta_5$ denominator and Mukai
lattice belong to the K3 row. The mixed
$T$-line/$J$-line profile is the diagnostic that the algebra is
non-principal DS (partition $(2,1)$); the global tag remains $\mathbf{M}$.
Witnesses: Fehily--Kawasetsu--Ridout~\cite{FKR20} for the OPE and
$c(k)$; Creutzig--Linshaw--Nakatsuka--Sato~\cite{CLNS24} for the
hook-type DS/bar transport sending $\mathcal{B}^k\mapsto\mathcal{B}^{-k-6}$.
\end{remark}


\section{Definition and OPE data}

\begin{definition}[Bershadsky--Polyakov algebra]
\label{def:bp-algebra}
\index{Bershadsky--Polyakov algebra}
\index{W-algebra!non-principal}
Let $k \in \mathbb{C} \setminus \{-3\}$. The
\emph{Bershadsky--Polyakov algebra}
$\mathcal{B}^k = \mathcal{W}^k(\mathfrak{sl}_3, f_{(2,1)})$ is the
BRST cohomology of the quantum Drinfeld--Sokolov reduction of
$V_k(\mathfrak{sl}_3)$ at the minimal nilpotent element
$f \in \mathfrak{sl}_3$ corresponding to the partition $(2,1)$
of~$3$. The Jacobson--Morozov $\mathfrak{sl}_2$-triple $(e,h,f)$
determines an $\mathrm{ad}(h)$-grading on $\mathfrak{sl}_3$; the
centralizer $\mathfrak{g}^f$ has dimension~$4$ (the number of strong
generators of $\mathcal{W}^k(\mathfrak{g}, f)$ equals $\dim \mathfrak{g}^f$
by Kac--Roan--Wakimoto), and the DS reduction produces four strong generators:
\begin{center}
\renewcommand{\arraystretch}{1.3}
\begin{tabular}{cccc}
\toprule
Generator & Conformal weight & Parity & $J$-charge \\
\midrule
$T$ & $2$ & bosonic & $0$ \\
$J$ & $1$ & bosonic & $0$ \\
$G^+$ & $3/2$ & bosonic/even & $+1$ \\
$G^-$ & $3/2$ & bosonic/even & $-1$ \\
\bottomrule
\end{tabular}
\end{center}
Here $J$ is the $\mathrm{U}(1)$ current from the residual Cartan line.
In the Feigin--Semikhatov normal form
\cite[(A.3)]{FeiginSemikhatov} its current level is
\[
k_J=\frac{2k+3}{3}.
\]
\end{definition}

\begin{proposition}[Feigin--Semikhatov OPE normal form for BP;
\ClaimStatusProvedElsewhere]
\label{prop:bp-ope-normal-form}
In the Feigin--Semikhatov normal form
\cite[(A.3)]{FeiginSemikhatov}, the bosonic OPE data
($T$, $J$ bosonic) are:
\begin{align}
T_{(3)}T &= \tfrac{c}{2}, &
T_{(1)}T &= 2T, &
T_{(0)}T &= \partial T,
\label{eq:bp-ope-TT} \\
J_{(1)}J &= k_J=\frac{2k+3}{3}, &
J_{(0)}J &= 0, &
T_{(1)}J &= J.
\label{eq:bp-ope-bosonic}
\end{align}
The conformal action on the charged even generators is
\begin{align}
T_{(1)}G^\pm &= \tfrac{3}{2}G^\pm, &
T_{(0)}G^\pm &= \partial G^\pm, &
J_{(0)}G^\pm &= \pm G^\pm .
\label{eq:bp-ope-primary}
\end{align}
The charged OPE data are
\begin{align}
G^+_{(2)}G^- &= (k+1)(2k+3), &
G^+_{(1)}G^- &= 3(k+1)J, &
G^+_{(0)}G^- &= 3{:}JJ{:}+\tfrac{3}{2}(k+1)\partial J-(k+3)T,
\label{eq:bp-ope-GpGm}
\end{align}
The reverse products are determined by ordinary vertex-algebra skew
symmetry; this is the even-field convention of
Fehily--Kawasetsu--Ridout~\cite[Definition~2.1]{FKR20}.
The self-OPEs $G^\pm_{(n)} G^\pm = 0$ vanish identically. The
maximal OPE pole order among generators is $p_{\max}(\mathcal{B}^k) = 4$
(from $T_{(3)}T$), and the collision depth is $k_{\max} = 3$.
\end{proposition}

\begin{convention}[BP working loci]
\label{conv:bp-generic-loci}
The universal Feigin--Semikhatov OPE presentation is used on
\[
U_{\mathrm{OPE}}=\mathbb{C}\setminus\{-3\}.
\]
The nondegenerate generator-pairing locus removes the residual-current
and charged-pairing zeros:
\[
U_{\mathrm{nd}}
=U_{\mathrm{OPE}}\setminus\{-1,-3/2\},
\qquad
k_J=\frac{2k+3}{3}\ne0,\quad
(k+1)(2k+3)\ne0.
\]
The Virasoro $T$-line shadow coefficients
$S_4^T$ and $S_5^T$ require
\[
U_T
=U_{\mathrm{OPE}}\setminus
\left\{
-\frac32,-\frac13,
\frac{-11\pm3\sqrt{89}}{20}
\right\},
\]
equivalently $c(k)\ne0$ and $5c(k)+22\ne0$. The sigma-invariant uses
the paired locus $U_T^\vee=\{k\in U_T\mid -k-6\in U_T\}$.
The curve-level chiral Hochschild statements use
$U_{\mathrm{Hoch}}$, the intersection of these loci with the PBW
Koszul, $E_\infty$-completed, finite-type/perfectness hypotheses of
Theorem~\ref{thm:hochschild-polynomial-growth}.  Admissible and
minimal-model simple quotient levels enter after a separate
simple-quotient comparison theorem.
\end{convention}


\section{Central charge}\label{sec:bp-central-charge}

\begin{proposition}[BP central charge; \ClaimStatusProvedHere]
\label{prop:bp-central-charge}
\index{Bershadsky--Polyakov algebra!central charge}
The central charge of $\mathcal{B}^k$ is
\begin{equation}\label{eq:bp-central-charge}
c(k) \;=\; -\frac{(2k+3)(3k+1)}{k+3}
\;=\; 25 - 6(k+3) - \frac{24}{k+3}\,.
\end{equation}
This is the Kac--Roan--Wakimoto normalization
used by Fehily--Kawasetsu--Ridout
\cite[Eq.~\textup{(}2.2\textup{)}]{FKR20}.
\end{proposition}

\begin{proof}
The quantum Drinfeld--Sokolov reduction of
$V_k(\mathfrak{sl}_3)$ at the minimal nilpotent has the
Feigin--Semikhatov strong-generator packet
\eqref{eq:bp-ope-bosonic}--\eqref{eq:bp-ope-GpGm}. The
Kac--Roan--Wakimoto formula for the minimal reduction
$\mathcal{W}^k(\fg, f_{\min})$ is
$c = k\dim\fg/(k+h^\vee) - 6k + h^\vee - 4$; for
$\fg = \mathfrak{sl}_3$ ($\dim\fg = 8$, $h^\vee = 3$) this is
\[
c(k)=\frac{8k}{k+3}-6k-1
=-\frac{6k^2+11k+3}{k+3}
=-\frac{(2k+3)(3k+1)}{k+3}.
\]
Two checks fix the normalization. At the admissible level
$k=-3/2$ the factor $2k+3$ vanishes and $c=0$. Under the
Feigin--Frenkel involution
$k\mapsto -k-6$ (that is, $k+3 \mapsto -(k+3)$), the odd part
$-6(k+3)-24/(k+3)$ cancels and direct addition gives
$c(k)+c(-k-6)=50$, the conductor used in
Theorem~\ref{thm:bp-koszul-conductor-polynomial}.
\end{proof}

\begin{remark}[Normalization]
\label{rem:bp-convention}
The level $k$ is the affine $\mathfrak{sl}_3$ level before
minimal Drinfeld--Sokolov reduction. The formula
\eqref{eq:bp-central-charge} is the one compatible with the
Feigin--Semikhatov OPE constants
\[
J_{(1)}J=\frac{2k+3}{3},\qquad
G^+_{(2)}G^-=(k+1)(2k+3),
\]
and with the conductor identity $c(k)+c(-k-6)=50$.
\end{remark}


\section{Verdier companion and conductor}\label{sec:bp-koszul}

The partition $(2,1)$ is self-transpose: $(2,1)^t = (2,1)$. The
objects remain distinct:
\[
A=\mathcal{B}^k,\qquad B(A),\qquad
A^{\mathrm i}=C_X(s^{-1}V,s^{-2}R),\qquad
K_X(A)=\mathbb D_{\Ran}B^{\mathrm{ch}}(A),\qquad
Z^{\mathrm{der}}_{\mathrm{ch}}(A).
\]
Bar--cobar reconstruction identifies $\Omega(B(A))$ with $A$.
The maps $q_A$, $\mathbb D(q_A)$, and
$\nu_A\colon K_X(A)\to A^!$ give the quadratic and chosen-pair
comparisons. Under the hook-type transport theorem
(Fehily~\cite{FehilyHook}, Creutzig--Linshaw--Nakatsuka--Sato
\cite{CLNS24}), this Verdier companion preserves the nilpotent orbit
and acts on the level via the Feigin--Frenkel involution of the parent
algebra $\widehat{\mathfrak{sl}}_3$.

\begin{theorem}[Bershadsky--Polyakov central-charge conductor identity; \ClaimStatusProvedHere]
\label{thm:bp-koszul-conductor-polynomial}
\index{Bershadsky--Polyakov algebra!central-charge conductor!polynomial identity}
\index{central-charge conductor!polynomial identity}
The calculation is made in the Open quadrant, in the
Feigin--Semikhatov chain-level normal form, at Beilinson level~$5$,
over the rational-function field $\mathbb{Q}(k)$, prior to admissible
or simple-quotient specialization.

In the rational-function field $\mathbb{Q}(k)$, the central-charge conductor
function
\[
K^c_{\mathcal{B}}(k)
\;:=\;
c(\mathcal{B}^k) + c(\mathcal{B}^{-k-6})
\;=\;
\Bigl(25 - 6(k+3) - \frac{24}{k+3}\Bigr)
 + \Bigl(25 + 6(k+3) + \frac{24}{k+3}\Bigr)
\]
is identically equal to $50$:
\[
K^c_{\mathcal{B}}(k) \;\equiv\; 50 \qquad \text{as an element of }
\mathbb{Q}(k).
\]
Equivalently, $c(\mathcal{B}^k) - 25$ is an odd function of $(k+3)$:
$[c(\mathcal{B}^k) - 25] + [c(\mathcal{B}^{-k-6}) - 25] = 0$ in
$\mathbb{Q}(k)$. The fixed point $k = -3$ of the level involution
$k \mapsto -k-6$ coincides with the critical level
$-h^\vee(\mathfrak{sl}_3) = -3$, where both summands are singular.
The polynomial identity holds globally, and the regular midpoint fibre
$c(\mathcal{B}^k)=25$ consists of $k=-3\pm2i$.
\end{theorem}

\begin{proof}
Direct polynomial-identity computation in $\mathbb{Q}(k)$. Write
$t := k+3$, so $c(\mathcal{B}^k) = 25 - 6t - 24/t$ and the level
involution $k \mapsto -k-6$ is $t \mapsto -t$. Then
\begin{align*}
K^c_{\mathcal{B}}(k)
&= \Bigl(25 - 6t - \frac{24}{t}\Bigr)
 + \Bigl(25 + 6t + \frac{24}{t}\Bigr)
 = 50\,.
\end{align*}
The level-dependent part $-6t - 24/t$ is odd in~$t$, so the
cancellation is an identity of rational functions. Hence $K^c_{\mathcal{B}}(k) - 50
\equiv 0$ in $\mathbb{Q}(k)$.
\end{proof}

\begin{proposition}[BP Verdier companion involution and conductor; \ClaimStatusConditional]
\label{prop:bp-self-duality}
\index{Bershadsky--Polyakov algebra!Verdier-Koszul companion}
\index{Koszul duality!non-principal Verdier companion}
The assertion lies in the Open quadrant, compares Beilinson
levels~$1$ and~$2$ by the quadratic comparison and pairwise Verdier
data of Theorems~A and~B,
and records the level~$5$ scalar shadow. The hypotheses are:
finite-type-or-completed bar-dual coalgebra, non-critical level
$k\neq-3$, and the subregular hook-type DS/bar transport statement
\cite{CLNS24,FehilyHook}.

For $k \neq -3$ \textup{(}away from the critical level $k = -h^\vee$\textup{)},
choose $A^{\mathrm i}=C_X(s^{-1}V,s^{-2}R)$ for
$A=\mathcal B^k$ and assume that
$q_A\colon A^{\mathrm i}\to B^{\mathrm{ch}}(A)$ is a finite-weight
quasi-isomorphism, or a filtered quasi-isomorphism in the completed
continuous Verdier pairing. Assume also the subregular DS/bar
transport statement.
Then the same-family Verdier-Koszul companion of $\mathcal{B}^k$ is
$(\mathcal{B}^k)^! \simeq \mathcal{B}^{k'}$ with
$k' = -k - 6$. The dual central charge is $c' = c(-k-6) = 50 - c$.
The central-charge conductor satisfies
\begin{equation}\label{eq:bp-conductor}
K^c_{\mathcal{B}} \;=\; c(k) + c(-k-6) \;=\; 50
\qquad \text{as a polynomial identity in } \mathbb{Q}(k)
\end{equation}
per Theorem~\ref{thm:bp-koszul-conductor-polynomial}.
\end{proposition}

\begin{proof}
The polynomial identity~\eqref{eq:bp-conductor} is
Theorem~\ref{thm:bp-koszul-conductor-polynomial}. The
self-transposition of $(2,1)$ and the Feigin--Frenkel involution
identify the companion level
$k^\vee = -k - 2h^\vee = -k-6$
\textup{(}where $h^\vee(\mathfrak{sl}_3) = 3$\textup{)}.
Under the assumed subregular DS/bar transport statement on the DS
quotient, the hook-transport theorem~\cite{CLNS24,FehilyHook} gives
$(\mathcal{B}^k)^! \simeq \mathcal{B}^{-k-6}$.
\end{proof}

\begin{remark}[Comparison of central-charge conductors]
\label{rem:bp-conductor-comparison}
\index{central-charge conductor!comparison}
The central-charge conductor grows rapidly with the complexity of the
DS reduction:
\begin{center}
\renewcommand{\arraystretch}{1.2}
\begin{tabular}{lcc}
\toprule
Algebra & Partition & $K$ \\
\midrule
$\mathrm{Vir}_c$ & (principal, $\mathfrak{sl}_2$) & $26$ \\
$\mathcal{W}_3$ (principal) & $(3)$ & $100$ \\
$\mathcal{B}^k$ (minimal) & $(2,1)$ & $50$ \\
\bottomrule
\end{tabular}
\end{center}
The principal $\mathcal{W}_3$ conductor $K = 100$ arises from a
different central charge parametrization: $c_{\mathcal{W}_3}(k) =
2 - 24(k{+}2)^2/(k{+}3)$. The self-dual point
$c_{\mathrm{sd}} = K/2$ is $c = 13$ for Virasoro, $c = 50$ for
$\mathcal{W}_3$, and $c = 25$ for the BP algebra.  The BP midpoint
$c=25$ lies on the complex level fibre: with $t=k+3$, the AM--GM bound
on $c = 25 - 6t - 24/t$ gives
$c \leq 1$ for $k > -3$ and $c \geq 49$ for
$k < -3$, leaving the interval $(1, 49) \ni 25$ unattained.
The formal self-dual point requires complex level
$k = -3 \pm 2i$.
\end{remark}

\begin{proposition}[Unified shift formula for $\mathfrak{sl}_3$
central-charge conductors; \ClaimStatusProvedHere]
\label{prop:sl3-conductor-shift-formula}
\index{central-charge conductor!sl3@$\mathfrak{sl}_3$!unified shift formula}
\index{shift parameter!nilpotent orbit}
For each non-zero nilpotent orbit
$\mathcal{O}_\lambda \subset \mathfrak{sl}_3$, let
$\mathcal{W}^k_\lambda := \mathcal{W}^k(\mathfrak{sl}_3, f_\lambda)$
denote the Drinfeld--Sokolov reduction at $f_\lambda \in \mathcal{O}_\lambda$.
Both non-trivial central charges admit the common parametrization
\begin{equation}\label{eq:sl3-conductor-c-shifted}
c\bigl(\mathcal{W}^k_\lambda\bigr)
\;=\;
c_{0,\lambda} \;-\; \frac{b_\lambda\bigl(k + h^\vee - s_\lambda\bigr)^2}{k + h^\vee},
\qquad h^\vee(\mathfrak{sl}_3) = 3,
\end{equation}
where the orbit-specific invariants
$(c_{0,\lambda}, b_\lambda, s_\lambda)$ are
\begin{equation}\label{eq:sl3-shift-table}
(c_{0}, b, s)_{(3)} \;=\; (2,\, 24,\, 1)\quad(\text{principal}),
\qquad
(c_{0}, b, s)_{(2,1)} \;=\; (1,\, 6,\, 2)\quad(\text{minimal}).
\end{equation}
Under the Feigin--Frenkel involution $k \mapsto k' = -k - 2h^\vee$,
the central-charge conductor evaluates to the closed-form Lie-theoretic
expression
\begin{equation}\label{eq:sl3-conductor-formula}
K_\lambda \;:=\; c\bigl(\mathcal{W}^k_\lambda\bigr)
+ c\bigl(\mathcal{W}^{k'}_\lambda\bigr)
\;=\; 2c_{0,\lambda} + 4\, b_\lambda s_\lambda
\quad \text{in } \mathbb{Q}(k).
\end{equation}
Specialising~\eqref{eq:sl3-conductor-formula} to the two non-trivial
$\mathfrak{sl}_3$ orbits recovers
$K_{(3)} = 4 + 96 = 100$ and $K_{(2,1)} = 2 + 48 = 50$.
\end{proposition}

\begin{proof}
The presentation~\eqref{eq:sl3-conductor-c-shifted} is read off from
the Kac--Roan--Wakimoto central-charge formula: for the principal
reduction \textup{(}Fateev--Lukyanov\textup{)},
$c = 50 - 24t - 24/t = 2 - 24(t-1)^2/t$ with $t = k + h^\vee$;
for the minimal reduction \textup{(}Bershadsky--Polyakov,
Proposition~\textup{\ref{prop:bp-central-charge}}\textup{)},
$c = 25 - 6t - 24/t = 1 - 6(t-2)^2/t$. Note that the two orbits
differ in the quartic coefficient $b_\lambda$ as well as the shift
$s_\lambda$: the minimal reduction carries the rank-one Virasoro
coefficient $b=6$, while the principal branch carries
$b=N(N^2{-}1)=24$.
The Feigin--Frenkel involution $k' = -k - 6$ sends
$k + h^\vee \mapsto -(k + h^\vee)$, hence
$k + h^\vee - s_\lambda \mapsto -(k + h^\vee + s_\lambda)$, so
$(k + h^\vee - s_\lambda)^2 \mapsto (k + h^\vee + s_\lambda)^2$,
while the denominator $k + h^\vee$ flips sign once.
Adding the two summands,
\begin{align*}
K_\lambda
&= 2c_{0,\lambda} \;-\; \frac{b_\lambda(k + h^\vee - s_\lambda)^2}{k + h^\vee}
 \;+\; \frac{b_\lambda(k + h^\vee + s_\lambda)^2}{k + h^\vee} \\
&= 2c_{0,\lambda} \;+\; \frac{b_\lambda\bigl[(k + h^\vee + s_\lambda)^2
 - (k + h^\vee - s_\lambda)^2\bigr]}{k + h^\vee}.
\end{align*}
Using the difference-of-squares identity
$(A + B)^2 - (A - B)^2 = 4AB$ with $A = k + h^\vee$ and $B = s_\lambda$,
the numerator becomes $4 s_\lambda (k + h^\vee)$, and the
$(k + h^\vee)$ factor cancels the denominator exactly once as a
rational-function identity.
This yields $K_\lambda = 2c_{0,\lambda} + 4 b_\lambda s_\lambda$ in
$\mathbb{Q}(k)$, with
the level-independence reduced to the algebraic cancellation.
\end{proof}

\begin{remark}[Three conductor computations]
\label{rem:sl3-conductor-verification}%
\index{central-charge conductor!multi-path computation}
Equation~\eqref{eq:sl3-conductor-formula} has three independent
numerical faces.
First, direct rational-function expansion at $\lambda = (2,1)$,
$(c_0, b, s) = (1, 6, 2)$, gives
$K = 2 + 48 = 50$, matching
Theorem~\textup{\ref{thm:bp-koszul-conductor-polynomial}} and
Proposition~\textup{\ref{prop:bp-self-duality}}.
At $\lambda = (3)$, $(c_0, b, s) = (2, 24, 1)$, it gives $K = 4 + 96 = 100$,
matching the Freudenthal--de Vries-style formula
$K_N = 2(N{-}1)(2N^2 + 2N + 1) = 100$ at $N = 3$
\textup{(}Remark~\textup{\ref{rem:koszul-conductor-explicit}})\textup{.}
Second, the self-dual central charge is $c_* = K_\lambda / 2$, so
$c_*((2,1)) = 25$ and $c_*((3)) = 50$, matching the values quoted in
Remark~\textup{\ref{rem:bp-conductor-comparison}} and
\textup{\ref{rem:self-dual-complementarity}}.
Third, the principal branch has the proved exponent sum
$\varrho_{(3)}=5/6$ and hence
$K^\kappa_{(3)}=(5/6)(100)=250/3$.  The minimal branch has even
charged generators $G^\pm$; its genus-$1$ scalar and modular
complementarity are the open obligations of
Propositions~\textup{\ref{prop:bp-kappa}} and
\textup{\ref{prop:bp-complementarity}}.
\end{remark}

\begin{remark}[Why the shift parameter $s_\lambda$ is the genuine
Lie-theoretic invariant]
\label{rem:shift-parameter-lie-theoretic}%
\index{shift parameter!Lie-theoretic interpretation}
The triple $(c_{0,\lambda}, b_\lambda, s_\lambda)$ in
Proposition~\textup{\ref{prop:sl3-conductor-shift-formula}} encodes
the orbit-dependent data that distinguish
one $\mathfrak{sl}_3$ DS reduction from another at fixed $h^\vee = 3$.
The two orbits carry separate quartic coefficients.  The principal
orbit carries the Fateev--Lukyanov coefficient
$b_{(3)} = N(N^2{-}1)|_{N=3} = 24$, while the minimal orbit carries
the rank-one Virasoro coefficient $b_{(2,1)} = 6$ --- the same
coefficient as the principal $\mathfrak{sl}_2$ reduction
$c(\mathrm{Vir}) = 1 - 6(t-1)^2/t$, of which the
Bershadsky--Polyakov shape $c = 1 - 6(t-2)^2/t$ is the $s = 2$
translate. The conductor
$K_\lambda = 2c_{0,\lambda} + 4 b_\lambda s_\lambda$ factors the
level-independent part as
\textup{(}twice the even constant\textup{)} plus
\textup{(}difference-of-squares factor $4$\textup{)} times
\textup{(}quartic coefficient times shift\textup{)}. Whether an
analogous three-parameter form covers
$\mathcal{W}^k(\mathfrak{sl}_N, f_\lambda)$ at \emph{all} non-principal
orbits is conjectural; here we record only the $\mathfrak{sl}_3$
specialisation as unconditional. The unified $\mathfrak{sl}_3$ form is
an \emph{$\mathfrak{sl}_3$-rank-specific} bridge identity.
\end{remark}


\section{Modular characteristic}\label{sec:bp-kappa}

\begin{proposition}[BP genus-$1$ modular-characteristic obligation;
\ClaimStatusOpen]
\label{prop:bp-kappa}
\index{Bershadsky--Polyakov algebra!modular characteristic}
\index{kappa@$\kappa$!Bershadsky--Polyakov}
For $k\in U_{\mathrm{OPE}}$, define the modular characteristic, when
the genus-$1$ curved chiral complex has been constructed, by
\begin{equation}\label{eq:bp-kappa}
F_1(\mathcal B^k)=\frac{\kappa_{\mathrm{BP}}(k)}{24}.
\end{equation}
Its evaluation requires the full genus-$1$ nonseparating curvature of
the BP chiral/BRST complex, including charged ghosts, neutral fields,
the conformal improvement, and mixed $J,G^\pm,T$ channels.  The current
status of $\kappa_{\mathrm{BP}}(k)$ is open.
\end{proposition}

\begin{proof}
Fehily--Kawasetsu--Ridout~\cite[Definition~2.1]{FKR20} define BP as an
ordinary bosonic vertex algebra.  The strong-generator packet is
\[
J_{h=1}^{\mathrm{even}},\qquad
G^+_{h=3/2}^{\mathrm{even}},\qquad
G^-_{h=3/2}^{\mathrm{even}},\qquad
T_{h=2}^{\mathrm{even}}.
\]
Accordingly, the signed reciprocal-weight expression formerly used in
this chapter evaluates to
\[
1+\frac{1}{3/2}+\frac{1}{3/2}+\frac{1}{2}
=1+\frac23+\frac23+\frac12=\frac{17}{6}.
\]
The principal-$W$ reciprocal-exponent theorem applies to principal DS
reduction.  BP is the minimal/subregular reduction, so its genus-$1$
coefficient is supplied by the direct curvature computation named in
the proposition.  The exact conformal data available before that
calculation are
\[
S_2^T=\frac c2,\qquad
S_2^J=\frac{2k+3}{3},\qquad
G^+_{(2)}G^-=(k+1)(2k+3),\qquad
\Delta^{\mathrm{DS}}_{\mathrm{BP}}(k)=-(6k+1),
\]
with the last identity obtained from
$c_{\mathrm{BP}}(k)-8k/(k+3)$.
\end{proof}

\begin{proposition}[Modular-characteristic complementarity obligation;
\ClaimStatusOpen]
\label{prop:bp-complementarity}
\index{Bershadsky--Polyakov algebra!complementarity}
\index{complementarity!Bershadsky--Polyakov}
For $k\ne-3$, define the BP modular-characteristic companion sum by
\begin{equation}\label{eq:bp-complementarity}
K^\kappa_{\mathrm{BP}}(k)
:=\kappa_{\mathrm{BP}}(k)+\kappa_{\mathrm{BP}}(-k-6).
\end{equation}
On $H_{\mathrm{BP}}^{\mathrm{DS/bar}}$, the two arguments correspond to
the same-family Verdier--Koszul pair.  Constancy and numerical
evaluation of~\eqref{eq:bp-complementarity} await the genus-$1$
computation of Proposition~\ref{prop:bp-kappa}.  The implication
\[
\kappa_{\mathrm{BP}}(k)=\frac{c_{\mathrm{BP}}(k)}6
\quad\Longrightarrow\quad
K^\kappa_{\mathrm{BP}}=\frac{25}{3}
\]
is a conditional algebraic consequence of $K^c_{\mathrm{BP}}=50$.
\end{proposition}

\begin{proof}
The definition is Equation~\eqref{eq:bp-complementarity}.
Proposition~\ref{prop:bp-self-duality} supplies the algebra-level
companion interpretation under $H_{\mathrm{BP}}^{\mathrm{DS/bar}}$.
Under the displayed additional hypothesis,
Theorem~\ref{thm:bp-koszul-conductor-polynomial} gives
\[
K^\kappa_{\mathrm{BP}}
=\frac{c(k)+c(-k-6)}6=\frac{50}{6}=\frac{25}{3}.
\]
\end{proof}

\begin{proposition}[Curve-level chiral Hochschild theorem for BP; \ClaimStatusConditional]
\label{prop:bp-curve-level-chirhoch}
\index{Bershadsky--Polyakov algebra!chiral Hochschild}
\index{chiral Hochschild cohomology!Bershadsky--Polyakov}
The assertion lies in the Open quadrant and passes from a complete
chain-level chart to the derived chiral centre by Theorem~H.  Let
$S_{\mathrm{BP}},S_{\mathrm{BP}}^!\subset\mathbb Z_{\geq0}$ be
finite.  On $U_{\mathrm{Hoch}}$ of
Convention~\ref{conv:bp-generic-loci}, assume the family data
$H_H(\mathcal B^k;S_{\mathrm{BP}})$ and
$H_H(\mathcal B^{-k-6};S_{\mathrm{BP}}^!)$, the chosen-pair comparison
$(\mathcal B^k)^!\simeq\mathcal B^{-k-6}$, and a perfect degree-two
chain pairing.  Simple quotients carry their own chart comparison.

Assume also the BP center condition
\[
Z_{\mathrm{ch}}(\mathcal{B}^k)=\mathbb{C}\mathbf{1},
\qquad
Z_{\mathrm{ch}}(\mathcal{B}^{-k-6})=\mathbb{C}\mathbf{1},
\]
and the BP derivation calculation
\[
\operatorname{OutDer}^{\mathrm{fix}}_{\mathrm{ch}}(\mathcal{B}^k)=0,
\]
where ``fixed'' means that the Feigin--Semikhatov OPE coefficients at
the chosen level $k$ are part of the product structure.  Then
\[
 \operatorname{Supp}\ChirHoch_X^\bullet(\mathcal B^k)
 \subseteq S_{\mathrm{BP}},
\]
and its Hilbert polynomial has the form
\[
P_{\mathcal{B}^k}(t)
=\sum_{n\ge0}\dim\ChirHoch_X^n(\mathcal{B}^k)t^n
=1+t^2+
 \sum_{n\in S_{\mathrm{BP}}\setminus\{0,1,2\}}
 d_nt^n,
\qquad
d_n:=\dim\ChirHoch_X^n(\mathcal B^k).
\]
In particular, $S_{\mathrm{BP}}\subseteq\{0,1,2\}$ gives
$P_{\mathcal B^k}(t)=1+t^2$.
Topological Hochschild theory, categorical Hochschild theory, the
product-ambient bulk, and a physical bulk occupy separate targets.
The Feigin--Semikhatov scalar package --- the $J$-level, the
$T$-central charge, and the mixed $G^+G^-$ pairing --- is degree-two
product metadata.  The linear coefficient is the fixed-product outer
derivation dimension.
\end{proposition}

\begin{proof}
Theorem~\ref{thm:main-koszul-hoch} gives the support inclusion, and
Theorem~\ref{thm:hochschild-polynomial-growth} expresses the Hilbert
polynomial as the sum over $S_{\mathrm{BP}}$.  The first center
hypothesis gives the constant coefficient~$1$.  The chosen-pair
comparison, the second center hypothesis, and the perfect degree-two
pairing give the quadratic coefficient~$1$.
The assumed fixed-product calculation
$\operatorname{OutDer}^{\mathrm{fix}}_{\mathrm{ch}}(\mathcal B^k)=0$
sets the linear coefficient to zero.  The remaining summands are
exactly the displayed degrees in
$S_{\mathrm{BP}}\setminus\{0,1,2\}$.
\end{proof}


\section{Shadow obstruction tower}\label{sec:bp-shadow}

The Bershadsky--Polyakov algebra has two bosonic generator lines:
the $T$-line (Virasoro restriction) and the $J$-line ($\mathrm{U}(1)$
restriction). The shadow obstruction tower is computed independently
on each line.

\subsection{The T-line: class M}

On the $T$-line, the shadow projections $S_r$ are determined
entirely by the Virasoro self-OPE $T_{(3)}T = c/2$, $T_{(1)}T = 2T$,
$T_{(0)}T = \partial T$. The shadow data coincides with that of
$\mathrm{Vir}_{c(k)}$:

\begin{proposition}[T-line shadow depth; \ClaimStatusConditional]
\label{prop:bp-tline-depth}
\index{Bershadsky--Polyakov algebra!shadow tower!T-line}
For $k\in U_T$, in the Virasoro single-line Riccati regime, the
shadow obstruction tower of $\mathcal{B}^k$
on the $T$-line has infinite depth:
$r_{\max}^T = \infty$. The tower is class~$M$. The shadow data is
\begin{equation}\label{eq:bp-tline-shadows}
S_2^T = \frac{c}{2}, \qquad
S_3^T = 2, \qquad
S_4^T = \frac{10}{c(5c+22)}, \qquad
S_5^T = \frac{-48}{c^2(5c+22)}, \qquad
\Delta_T = \frac{40}{5c+22}\,.
\end{equation}
Equivalently, after substituting
$c=c_{\mathrm{BP}}(k)=-(2k+3)(3k+1)/(k+3)$,
\[
S_4^T
=\frac{10(k+3)^2}
{3(2k+3)(3k+1)(10k^2+11k-17)},\qquad
S_5^T
=\frac{16(k+3)^3}
{(2k+3)^2(3k+1)^2(10k^2+11k-17)}.
\]
The per-line modular characteristic is $\kappa_T = c/2$
\textup{(}the Virasoro value, distinct from the full-algebra
$\kappa_{\mathrm{BP}}(k)$ of Proposition~\ref{prop:bp-kappa}\textup{)}.
\end{proposition}

\begin{proof}
On the $T$-line, the MC equation projects to the single-line
recursion with $\kappa_T = c/2$, $\alpha = S_3 = 2$, and
$S_4 = Q^{\mathrm{contact}}_{\mathrm{Vir}} = 10/(c(5c+22))$.
The next Virasoro coefficient is $S_5=-48/[c^2(5c+22)]$, and direct
substitution of $c_{\mathrm{BP}}(k)$ gives the displayed rational form.
The critical discriminant is $\Delta_T = 8 \kappa_T S_4 =
8 \cdot (c/2) \cdot 10/(c(5c+22)) = 40/(5c+22)$. Since
$k\in U_T$ gives $c\ne0$ and $5c + 22 \neq 0$,
$\Delta_T \neq 0$ and the shadow metric is irreducible.
The tower has infinite depth.
\end{proof}

\begin{remark}[Quartic line separation]
\label{rem:bp-quartic-line-separation}
The coefficient $S_4^T$ in
Proposition~\ref{prop:bp-tline-depth} is the Virasoro
$T$-line quartic contact. A scalar quartic coefficient on the
mixed charged bilinear line generated by $G^+G^-$ is a different
invariant: the Feigin--Semikhatov simple pole contains
$3{:}JJ{:}+\tfrac32(k+1)\partial J-(k+3)T$.  The mixed line therefore
requires its own projection.
\end{remark}


\subsection{The J-line: class G}

The $J$-line is spanned by the $\mathrm{U}(1)$ current. The
$J$-self-OPE has a single double pole $J_{(1)}J = (2k+3)/3$
and $J_{(0)}J = 0$: a purely quadratic
(Heisenberg-type) structure at the residual level.

\begin{proposition}[J-line shadow depth; \ClaimStatusProvedHere]
\label{prop:bp-jline-depth}
\index{Bershadsky--Polyakov algebra!shadow tower!J-line}
For $k\in U_{\mathrm{OPE}}$, the $J$-line shadow coefficients are
\[
S_2^J=\kappa_J=\frac{2k+3}{3},\qquad
S_r^J=0\quad(r\ge3).
\]
If $k\ne -3/2$, this is a nondegenerate class~$G$ tower with
$r_{\max}^J=2$. At $k=-3/2$ the residual current level vanishes and
the $J$-line becomes a degenerate Gaussian line.
\end{proposition}

\begin{proof}
The $J$-self-OPE is purely quadratic: $J_{(1)}J = (2k+3)/3$
and $J_{(0)}J = 0$. The cubic shadow requires a composite-field
contribution from $J_{(0)}J$, which vanishes. Higher degrees
vanish by induction: the shadow obstruction tower terminates at
degree~$2$.
\end{proof}

\begin{remark}[Mixed depth: the first non-principal witness]
\label{rem:bp-mixed-depth}
\index{shadow depth!mixed}
For $k\in U_{\mathrm{nd}}\cap U_T$, the Bershadsky--Polyakov algebra
has shadow depth $\infty$ on the $T$-line and shadow depth~$2$ on the
$J$-line. In the principal
landscape, each algebra has a single shadow class
(Heisenberg is~$G$ on every line, $\widehat{\mathfrak{sl}}_N$
is~$L$, Virasoro is~$M$). The BP algebra is the first
non-principal example of mixed shadow depth within a single algebra.
The global shadow class is $M$ (the maximum over all lines).
The per-line classification is
$(T\colon M,\; J\colon G,\; G^\pm\colon \text{even charged})$.
The $\mathcal{N}=2$ SCA of Chapter~\ref{chap:n2-sca} has the same
conformal weights and the opposite parity on $G^\pm$; the two OPE
algebras therefore occupy different superselection types.
\end{remark}


\subsection{Sigma-invariant}

The sigma-invariant measures the complementarity defect at each
degree:
\begin{equation}\label{eq:bp-sigma}
\sigma^{(r)}(k) \;:=\; S_r^T(c(k)) + S_r^T(c(-k-6)),
\qquad c' = 50 - c\,.
\end{equation}

\begin{proposition}[The first three sigma functions; \ClaimStatusConditional]
\label{prop:bp-sigma}
\index{Bershadsky--Polyakov algebra!sigma-invariant}
For $k\in U_T^\vee$, the $T$-line sigma functions satisfy
$\sigma^{(2)} = 25$,
$\sigma^{(3)} = 4$,
\[
\sigma^{(4)}
=
\frac{10}{c(5c+22)}
+
\frac{10}{(50-c)(272-5c)},
\]
which is a nonzero rational function of~$k$.  Higher sigma functions
are defined by Equation~\eqref{eq:bp-sigma}; their zero divisors require
degree-by-degree computation.
\end{proposition}

\begin{proof}
At degree~$2$: $\sigma^{(2)} = c/2 + (50-c)/2 = 25$.
At degree~$3$: $\sigma^{(3)} = 2 + 2 = 4$.
At degree~$4$, substitute $c'=50-c$ into the Virasoro quartic
contact coefficient $S_4=10/[c(5c+22)]$:
\[
\sigma^{(4)}
=\frac{10}{c(5c+22)}
 +\frac{10}{c'(5c'+22)}
=\frac{10}{c(5c+22)}
 +\frac{10}{(50-c)(272-5c)}.
\]
The numerator of this rational function is nonzero in~$\mathbb Q(c)$.
Its polar divisors in
the $c$-line are $c=0$, $c=50$, $c=-22/5$, and $c=272/5$. After
$c=c_{\mathrm{BP}}(k)$ the pole pairs are
$(c,c')=(0,50)$ and $(-22/5,272/5)$, matching
Computation~\ref{comp:bp-shadow-tower}.  This proves the stated
degrees~$2$, $3$, and~$4$.
\end{proof}


\section{Comparison with $\mathcal{W}_3$ and the
$\mathcal{N}=2$ SCA}\label{sec:bp-comparison}

The three DS reductions of $\widehat{\mathfrak{sl}}_3$ are the
full affine algebra (trivial orbit), the Bershadsky--Polyakov
algebra (minimal orbit), and $\mathcal{W}_3$ (principal orbit).
The comparison:

\begin{center}
\renewcommand{\arraystretch}{1.3}
\begin{tabular}{l|ccc}
\toprule
& $\widehat{\mathfrak{sl}}_{3,k}$ &
$\mathcal{B}^k$ (minimal) &
$\mathcal{W}_3$ (principal) \\
\midrule
Partition & $(1,1,1)$ & $(2,1)$ & $(3)$ \\
Generators & $J^a$ ($a = 1,\ldots,8$) & $T, J, G^\pm$ & $T, W$ \\
$\#$ generators & $8$ & $4$ & $2$ \\
Generator weights & $1$ & $2, 1, 3/2, 3/2$ & $2, 3$ \\
$\kappa$ & $\dfrac{4(k+3)}{3}$ & open genus-$1$ value & $\dfrac{5c}{6}$ \\
$\varrho = \kappa/c$ & $\dfrac{(k+3)^2}{6k}$ & open & $\dfrac{5}{6}$ \\
$K$ & $16$ & $50$ & $100$ \\
Shadow class & L & M (mixed) & M \\
\bottomrule
\end{tabular}
\end{center}

The row $\varrho=\kappa/c$ records an instantaneous ratio after
$\kappa$ has been computed.  For the principal $\mathcal W_3$ branch,
the exponent theorem gives $\varrho=5/6$.  For
$\widehat{\mathfrak{sl}}_{3,k}$ the instantaneous ratio
$(k+3)^2/(6k)$ is level-dependent, while the companion sum of the
affine modular characteristics is~$0$.  For the minimal BP branch, the
direct genus-$1$ computation of Proposition~\ref{prop:bp-kappa} is the
remaining input.

\begin{remark}[Relation to the $\mathcal{N}=2$ SCA]
\label{rem:bp-n2-comparison}
\index{Bershadsky--Polyakov algebra!relation to N=2 SCA}
The BP algebra and the $\mathcal{N}=2$ SCA share the conformal-weight
profile $(2,1,3/2,3/2)$ and the per-line shadow labels
($T\colon M$, $J\colon G$).  Their parity data separate them:
$G^\pm$ are even in BP and odd in the $\mathcal N=2$ SCA.  Their OPE
structures also differ:
the $\mathcal{N}=2$ SCA has
$J_{(1)}J = c/3$ and $G^+_{(2)}G^- = c/3$
(Chapter~\ref{chap:n2-sca}, eq.~\eqref{eq:n2-ope-JJ}
and~\eqref{eq:n2-ope-GG}),
while the BP algebra has
\[
J_{(1)}J=\frac{2k+3}{3},\qquad
G^+_{(2)}G^-=(k+1)(2k+3),
\]
and the simple pole of $G^+G^-$ contains
$3{:}JJ{:}+\tfrac32(k+1)\partial J-(k+3)T$.  The
$\mathcal N=2$ modular characteristic has its own family formula;
$\kappa_{\mathrm{BP}}(k)$ is the open genus-$1$ curvature coefficient
of Proposition~\ref{prop:bp-kappa}.
\end{remark}


\section{The hook-type family}\label{sec:bp-hook}

The Bershadsky--Polyakov algebra is the seed of the hook-type
transport family. A hook partition
$\lambda = (N-j, 1^j)$ of $N$ is self-transpose if and only if
$N$ is odd and $j = (N-1)/2$
(Proposition~\ref{prop:bp-hook-series} below).

\begin{proposition}[Self-transpose hook companion involution; \ClaimStatusConditional]
\label{prop:bp-hook-series}
\index{hook-type family}
For every odd $N = 2m+1$, at non-critical level
($k \neq -h^\vee = -N$), and $\lambda = (m+1, 1^m)$, assume the
hook-type DS/bar transport statement and a finite-type or completed
continuous Verdier pairing on the bar-dual coalgebra. Then the
$\mathcal{W}$-algebra
$\mathcal{W}_k(\mathfrak{sl}_N, f_\lambda)$ has same-family
Verdier-Koszul companion
$(\mathcal{W}_k(\mathfrak{sl}_N, f_\lambda))^!
\simeq
\mathcal{W}_{k'}(\mathfrak{sl}_N, f_\lambda)$
with $k' = -k - 2N$ \textup{(}where $h^\vee(\mathfrak{sl}_N) = N$,
so $k' = -k - 2h^\vee$\textup{)}.
The first three cases are:
\begin{center}
\renewcommand{\arraystretch}{1.2}
\begin{tabular}{cccc}
\toprule
$N$ & $\lambda$ & Algebra & Dual level shift \\
\midrule
$3$ & $(2,1)$ & $\mathcal{B}^k$ & $k' = -k - 6$ \\
$5$ & $(3,1,1)$ & $\mathcal{W}^k(\mathfrak{sl}_5, f_{(3,1,1)})$
 & $k' = -k - 10$ \\
$7$ & $(4,1,1,1)$ & $\mathcal{W}^k(\mathfrak{sl}_7, f_{(4,1,1,1)})$
 & $k' = -k - 14$ \\
\bottomrule
\end{tabular}
\end{center}
\end{proposition}

\begin{proof}
The hook partition $(N-j,1^j)$ is self-transpose precisely when its
first row and first column have the same length, hence when
$N-j=j+1$, i.e. $N=2j+1$. Thus for $N=2m+1$ the self-transpose hook is
$(m+1,1^m)$. The Feigin--Frenkel involution of
$\widehat{\mathfrak{sl}}_N$ sends $k$ to $-k-2h^\vee=-k-2N$.
Under the assumed hook-type DS/bar transport and the Verdier
finite-type or completion hypothesis, the bar-dual companion remains
in the same nilpotent orbit, so the displayed same-family
identification follows.
\end{proof}

\begin{corollary}[Hook-series seed: BP conductor; \ClaimStatusConditional]
\label{cor:bp-koszul-conductor-seed}
The Bershadsky--Polyakov algebra
$\mathcal{B}^k = \mathcal{W}_k(\mathfrak{sl}_3, f_{(2,1)})$
is the $N=3$ case of Proposition~\ref{prop:bp-hook-series}: at
non-critical level $k \neq -3$, the central-charge conductor identity
\[
 K^c_{\mathcal{B}}(k) \;:=\;
 c(\mathcal{B}^k) + c(\mathcal{B}^{-k-6}) \;\equiv\; 50
 \qquad \text{in } \mathbb{Q}(k)
\]
of Theorem~\ref{thm:bp-koszul-conductor-polynomial} seeds the
hook-series: every self-transpose hook of odd $N$ produces a
same-family companion fixed locus, with algebraic Verdier-Koszul
companionhood conditional on hook-type DS/bar transport and the
finite-type or completed Verdier pairing; BP ($N=3$) is the first
entry.
\end{corollary}

\begin{proof}
Set $N=3$ and $m=1$ in
Proposition~\ref{prop:bp-hook-series}. The self-transpose hook is
$(2,1)$, $h^\vee(\mathfrak{sl}_3)=3$, and the companion level is
$-k-6$. The conductor identity is exactly
Theorem~\ref{thm:bp-koszul-conductor-polynomial}.
\end{proof}

\begin{remark}[Bala--Carter dependence of the conductor]
\label{rem:kw-bp-vs-principal}
The Bershadsky--Polyakov conductor $K^c_{\mathcal{B}} = 50$
differs from the $\mathcal{W}_3$ conductor $K^c_{\mathcal{W}_3} = 100$
\textup{(}principal $\mathfrak{sl}_3$, partition $(3)$\textup{)},
reflecting the Bala--Carter dependence
$(\mathfrak{g}, f) \mapsto K^c$. The formula
$K = 2r + 12\sum_i m_i(m_i+1)$ of the W-algebra chapter has the
principal nilpotent as its stated domain. BP and $\mathcal{W}_3$ have
distinct same-family companion involutions ($k \mapsto -k-6$ for BP
versus the principal involution for $\mathcal{W}_3$); interpreting
these fixed loci as algebraic Verdier-Koszul companion loci is
conditional on the corresponding DS/bar transport statement and
Verdier finiteness/completion.
\end{remark}
